\chapter{Συζήτηση, Περιορισμοί και Μελλοντική Εργασία}
\label{ch:discussion}

\section{Σύνοψη ευρημάτων}
\label{sec:disc-summary}

Τα αποτελέσματα του Κεφαλαίου~\ref{ch:results} απαντούν στα τρία
ερευνητικά ερωτήματα της Ενότητας~\ref{subsec:intro-questions}:

\begin{itemize}
  \item[\textbf{Ε1.}] Η επίγεια και η δορυφορική εξυπηρέτηση διαφέρουν
    ριζικά και συστηματικά ως προς throughput, ενέργεια ανά bit, και SNR
    (Ενότητα~\ref{sec:results-repeated}): ο ρυθμός μετάδοσης και η
    ενεργειακή απόδοση ανά bit κινούνται σε αντίθετη κατεύθυνση μεταξύ των
    δύο τύπων ζεύξης, μια ισορροπία που κάθε μελλοντική πολιτική
    επιλογής/εναλλαγής κόμβου θα πρέπει να λαμβάνει υπόψη αν στοχεύει σε
    ενεργειακά αποδοτική λειτουργία.
  \item[\textbf{Ε2.}] Η χρονική προσομοίωση διέλευσης
    (Ενότητα~\ref{sec:results-temporal}) αναπαράγει με φυσική συνέπεια το
    αναμενόμενο σχήμα γωνίας ανύψωσης, και δείχνει ότι όσο ο δορυφόρος
    πλησιάζει το ζενίθ, το SNR/throughput του δορυφορικού σκέλους
    βελτιώνεται σταθερά, παραμένει όμως, ακόμα και στο ζενίθ, κατά πολύ
    χαμηλότερο από ένα καλό επίγειο link. Η δορυφορική ζεύξη δεν
    "νικά" ποτέ σε throughput στο συγκεκριμένο σενάριο, παρά μόνο σε
    ενεργειακή απόδοση και σε διαθεσιμότητα, εκεί όπου δεν υπάρχει καθόλου
    εναλλακτική.
  \item[\textbf{Ε3.}] Εντοπίστηκαν πέντε συγκεκριμένοι περιορισμοί του
    αρχικού μοντέλου (Ενότητες~\ref{subsec:results-handovers},
    \ref{subsec:results-outage}, \ref{subsec:results-hysteresis},
    \ref{subsec:results-fairness}, Ενότητα~\ref{sec:meth-node-selection}):
    μη-ρεαλιστικά συχνά handovers λόγω μη-συσχετισμένου shadow fading,
    απουσία ρητής διαχείρισης outage, αδυναμία μοντελοποίησης πραγματικής
    πολυσυνδεσιμότητας (ο αρχικός κανόνας ήταν καθαρά single-connectivity),
    απουσία υστέρησης/time-to-trigger στην ενεργοποίηση ζεύξης (καθαρά
    greedy κανόνας στιγμιαίου SNR, ευάλωτος σε τεχνητά συχνή εναλλαγή σε
    θορυβώδεις συνθήκες κοντά στο κατώφλι), και απουσία συντονισμού
    μεταξύ χρηστών στην απόφαση σύνδεσης (κάθε $\mathcal{N}(u)$
    υπολογιζόταν πλήρως ανεξάρτητα, χωρίς να λαμβάνεται υπόψη το φορτίο
    του κόμβου ή το αν άλλοι χρήστες του ίδιου κόμβου έχουν εναλλακτική).
    Και οι πέντε αντιμετωπίστηκαν ήδη εντός του Μέρους~1: ο πρώτος με τη
    χωρικά συσχετισμένη διαδικασία shadow fading της
    Ενότητας~\ref{subsec:meth-correlated-sf} (ο
    Πίνακας~\ref{tab:handovers} επιβεβαιώνει ότι τα handovers μειώθηκαν
    στο φυσικά αναμενόμενο επίπεδο -- πλέον μηδέν πραγματικά handovers σε
    ολόκληρη τη διέλευση), ο δεύτερος με το ελάχιστο χρησιμοποιήσιμο SNR
    κατώφλι της εξίσωσης~\eqref{eq:node-selection}
    (Ενότητα~\ref{subsec:results-outage} επιβεβαιώνει ότι ο χρήστης $5$
    καταγράφεται πλέον σε outage αντί να αναγκάζεται σε μια ζεύξη
    μηδενικής πρακτικής χρησιμότητας), ο τρίτος με το τετρα-καταστασιακό
    μοντέλο σύνδεσης SS-SBS της ίδιας εξίσωσης (ταυτόχρονη εξυπηρέτηση
    BS+δορυφόρου όποτε και οι δύο ξεπερνούν το κατώφλι), ο τέταρτος
    με μια μηχανή υστέρησης+TTT κατά το πρότυπο Event~A3 του TS~38.331
    (εξίσωση~\eqref{eq:hysteresis}, Ενότητα~\ref{subsec:meth-hysteresis}) --
    επικυρωμένη σε ειδικά κατασκευασμένο σενάριο πίεσης, αφού το ίδιο το
    σενάριο αναφοράς αποδείχθηκε πολύ "ήσυχο" (καθόλου ταλάντωση κοντά
    στο κατώφλι) για να την αναδείξει (Ενότητα~\ref{subsec:results-hysteresis}):
    $67\%$ μείωση μεταβάσεων ($12\rightarrow4$) σε $80$ βήματα -- και ο
    πέμπτος με μια load-based φάση εξισορρόπησης
    (Ενότητα~\ref{subsec:meth-fairness}) σε δύο αλγορίθμους: \emph{PerBs}
    (εξίσωση~\eqref{eq:fairness}), που αποσυνδέει προτεραιοποιημένα από
    έναν υπερφορτωμένο BS τους χρήστες που έχουν ήδη δορυφορική
    εναλλακτική, προστατεύοντας όσους δεν έχουν -- επικυρωμένος σε
    αντίστοιχο σενάριο συμφόρησης (Ενότητα~\ref{subsec:results-fairness}):
    $+33\%$ μέση/ελάχιστη χωρητικότητα για την προστατευόμενη ομάδα, με
    πραγματικό αλλά περιορισμένο κόστος για την ομάδα με εναλλακτική· και
    \emph{Joint} (εξίσωση~\eqref{eq:fairness-joint}), που προτιμά lateral
    handover σε γείτονα BS αντί για πτώση σε δορυφόρο-μόνο, και μπορεί να
    ωφελήσει ακόμα και Terrestrial-only χρήστες -- επικυρωμένος σε
    δεύτερο, δύο-BS σενάριο: μηδενική ανακούφιση από το PerBs έναντι
    πλήρους αναδιανομής φορτίου από το Joint σε μια Terrestrial-only
    παραλλαγή, με το ειλικρινές εύρημα ότι η μέση χωρητικότητα του
    πληθυσμού μπορεί να \emph{μειωθεί} όταν το Joint γεμίζει έναν
    προηγουμένως ελαφρά φορτωμένο γείτονα (Ενότητα~\ref{subsec:results-fairness}).
    Ως προς το δεύτερο
    σκέλος του
    ερωτήματος -- τι θα χρειαστεί ώστε ένα μελλοντικό μοντέλο μηχανικής
    μάθησης να αντιμετωπίσει αυτούς τους περιορισμούς -- η
    Ενότητα~\ref{sec:results-ml-realistic} δίνει μια πρώτη, μερική
    απάντηση, με μια ενδιάμεση παράκαμψη στην πορεία: η πολυσυνδεσιμότητα
    άλλαξε αρχικά την ίδια την κατανομή κλάσεων της ετικέτας ώστε η
    επίδοση με ρεαλιστικά διαθέσιμο (θορυβώδες) SNR να μείνει σχεδόν
    αμετάβλητη σε σχέση με το ιδανικό SNR ($99.15\%$ έναντι $99.72\%$
    Accuracy, Random Forest) -- όχι επειδή το πρόβλημα έγινε γενικά
    εύκολο, αλλά επειδή η ίδια η γεννήτρια σεναρίων δεν παρήγε πια
    οριακές περιπτώσεις. Μια επακόλουθη διόρθωση της γεννήτριας
    (Ενότητα~\ref{sec:results-ml-realistic}) αποκατέστησε ένα γνήσιο,
    μετρήσιμο χάσμα ($100.00\%$ ακριβές SNR έναντι $93.93\%$
    μόνο-γεωμετρίας, Random Forest). Και τα δύο ευρήματα συζητούνται
    αναλυτικά παρακάτω.
\end{itemize}

\section{Περιορισμοί του τρέχοντος μοντέλου}
\label{sec:disc-limitations}

Πέρα από τα δύο ευρήματα που αναδείχθηκαν εμπειρικά στο
Κεφάλαιο~\ref{ch:results}, το τρέχον μοντέλο έχει και άλλους γνωστούς
περιορισμούς, οι οποίοι δεν έχουν ακόμα αντιμετωπιστεί:

\begin{itemize}
  \item \textbf{Ακόμα greedy heuristic, όχι πλήρης joint
    βελτιστοποίηση ούτε πιο σύνθετος διαμοιρασμός πόρων.} Ο αλγόριθμος
    Joint της Ενότητας~\ref{subsec:meth-fairness} συντονίζει πλέον
    πολλαπλούς BS (δεν αντιμετωπίζει πια κάθε BS εντελώς μεμονωμένα, όπως
    το PerBs), αλλά παραμένει ένα greedy, per-iteration heuristic:
    διαλέγει προορισμό με βάση το ποιος BS έχει τη \emph{μεγαλύτερη
    ελεύθερη χωρητικότητα}, όχι με βάση μια αντικειμενική συνάρτηση πάνω
    στη συνολική/βεβαρημένη χωρητικότητα του δικτύου (π.χ. proportional
    fairness πάνω στο throughput) -- κάτι που, όπως δείχνει το εύρημα της
    Ενότητας~\ref{subsec:results-fairness}, μπορεί να μειώσει τη μέση
    χωρητικότητα του πληθυσμού ακόμα κι όταν κάθε μεμονωμένος
    μετακινημένος χρήστης βελτιώνεται. Ούτε μοντελοποιεί πιο σύνθετα
    σχήματα διαμοιρασμού (π.χ. Mode~2/3 της SS-SBS ή η αρχιτεκτονική
    MS-MBS, Ενότητα~\ref{subsec:bg-architectures}, με κοινή διαχείριση
    φάσματος/ισχύος μεταξύ πολλαπλών δορυφόρων ή σταθμών βάσης).
  \item \textbf{Δορυφορικό κανάλι χωρίς ιονοσφαιρική σπινθηρίδα.} Το
    μοντέλο FSPL περιλαμβάνει πλέον ατμοσφαιρική απόσβεση αερίων
    (εξίσωση~\eqref{eq:gas-atten}, TR~38.811 §6.6.4), αλλά όχι
    ιονοσφαιρική σπινθηρίδα (§6.6.6) -- το σημαντικότερο εναπομείναν
    φαινόμενο στη ζώνη S, το οποίο όμως απαιτεί κλιματολογικό μοντέλο
    (γεωγρ. πλάτος/ώρα/εποχή/ηλιακή δραστηριότητα) και όχι κλειστή
    αναλυτική σχέση. Ούτε κέρδος κεραίας εξαρτώμενο από τη γωνία
    (TR~38.901 §7.3) μοντελοποιείται σε καμία από τις δύο πλευρές: το
    επίγειο και το δορυφορικό σκέλος χρησιμοποιούν πλέον συμμετρικά ένα
    ενιαίο, βαθμωτό κέρδος κεραίας το καθένα
    (εξισώσεις~\eqref{eq:snr-bs},~\eqref{eq:snr-sat}), όχι όμως πλήρες
    μοτίβο ακτινοβολίας/πίνακα κεραιών με beamforming.
  \item \textbf{Απλοποιημένο μοντέλο ίχνους εδάφους.} Το μοντέλο κίνησης
    δορυφόρου (Ενότητα~\ref{sec:meth-temporal}) χρησιμοποιεί επίπεδη
    γεωμετρία σε σταθερό γεωγραφικό πλάτος, όχι πλήρη ορβιτογράφο (π.χ.
    SGP4 με πραγματικά two-line elements).
  \item \textbf{Ενεργειακό μοντέλο μόνο ενεργών κόμβων.} Η
    εξίσωση~\eqref{eq:earth-power} δεν συνυπολογίζει την κατανάλωση
    ισχύος $P_{\text{sleep}}$ κόμβων χωρίς εξυπηρετούμενους χρήστες,
    παρότι η παράμετρος αυτή ήδη υπάρχει στο μοντέλο EARTH.
\end{itemize}

\section{Μελλοντική εργασία}
\label{sec:disc-future-work}

Με βάση τα παραπάνω, και τις οδηγίες που δόθηκαν από τον επιβλέποντα
καθηγητή σχετικά με την ανάγκη διεξαγωγής περισσότερων, χρονικά εκτενών
προσομοιώσεων και την αξιολόγηση της συμπεριφοράς του μοντέλου, η εργασία
θα συνεχιστεί προς τις εξής κατευθύνσεις, σε φθίνουσα σειρά προτεραιότητας.

Οκτώ από τις κατευθύνσεις που αρχικά τέθηκαν εδώ ως μελλοντική
εργασία υλοποιήθηκαν ήδη εντός του Μέρους~1, και αναφέρονται εδώ ως
ολοκληρωμένες αντί να επαναλαμβάνονται στην αριθμημένη λίστα παρακάτω:

\begin{itemize}
  \item Ένα χρονικά συσχετισμένο μοντέλο shadow fading αντί για
    ανεξάρτητο δείγμα (εξίσωση~\eqref{eq:pathloss-shadowing}) ανά χρονικό
    βήμα (Ενότητα~\ref{subsec:meth-correlated-sf},
    εξίσωση~\eqref{eq:correlated-sf}), μέσω εκθετικής αυτοσυσχέτισης κατά
    Gudmundson~\cite{gudmundson1991correlation}, παραμετροποιημένης με
    την correlation distance του TR~38.901, Πίνακας~7.5-6. Τα
    αποτελέσματα (Πίνακας~\ref{tab:handovers}) επιβεβαιώνουν ότι τα
    handovers μειώθηκαν στο φυσικά αναμενόμενο επίπεδο.
  \item Φραγμένη χωρητικότητα MCS/CQI αντί για μη-φραγμένη χωρητικότητα
    Shannon (εξίσωση~\eqref{eq:shannon-capacity}, TS~38.214~\cite{ts38214}
    Πίνακας 5.1.3.1-2), η οποία υπερεκτιμούσε σημαντικά τη χωρητικότητα
    σε υψηλό SNR (βλ. Ενότητα~\ref{sec:results-static}).
  \item Κέρδος κεραίας στο επίγειο σκέλος (εξίσωση~\eqref{eq:snr-bs},
    TR~38.901 §7.3, Πίνακας 7.3-1) και ατμοσφαιρική απόσβεση αερίων στο
    δορυφορικό σκέλος (εξίσωση~\eqref{eq:gas-atten}, TR~38.811 §6.6.4),
    ώστε τα δύο σκέλη να βρίσκονται σε πιο συγκρίσιμο επίπεδο
    μοντελοποίησης· βλ. Ενότητα~\ref{sec:disc-limitations} για τι
    παραμένει ασύμμετρο (πλήρες μοτίβο ακτινοβολίας κεραίας, ιονοσφαιρική
    σπινθηρίδα).
  \item Πρόβλεψη από θορυβώδη ή ελλιπή παρατηρήσιμα αντί για το ακριβές,
    στιγμιαίο SNR και των δύο υποψηφίων (Ενότητα~\ref{sec:results-ml-realistic},
    Πίνακας~\ref{tab:ml-results-realistic}): δύο νέα πειράματα
    (\texttt{train\_model\_noisy\_snr.py}, \texttt{train\_model\_geometry\_only.py})
    έδειξαν αρχικά, αμέσως μετά την προσθήκη της πολυσυνδεσιμότητας (βλ.
    τελευταία κουκκίδα παρακάτω), σχεδόν ταυτόσημη Accuracy με το ακριβές
    SNR ($99.72\% \rightarrow 99.15\% \rightarrow 99.15\%$ για το Random
    Forest) αντί για τη σαφή, τρίβημη πτώση που παρατηρούνταν πριν --
    ένα εύρημα που αποδείχθηκε τεχνούργημα της τότε γεννήτριας σεναρίων
    (βλ. Ενότητα~\ref{subsec:disc-mc-ml-impact}): μετά τη διόρθωσή της, η
    πτώση επανεμφανίστηκε ($100.00\% \rightarrow 93.12\% \rightarrow
    93.93\%$).
  \item Ρητή κατάσταση outage (εξίσωση~\eqref{eq:node-selection},
    Ενότητα~\ref{subsec:results-outage}): ένα ελάχιστο χρησιμοποιήσιμο SNR
    κατώφλι $\text{SNR}_{\min} \approx -7.53$~dB, παραγόμενο από το ίδιο
    TS~38.214 Πίνακα 5.1.3.1-2 (MCS~0) που ήδη χρησιμοποιείται για το όριο
    χωρητικότητας παραπάνω. Όταν κανένας υποψήφιος δεν το ξεπερνά, ο
    χρήστης καταγράφεται ρητά ως εκτός κάλυψης (\texttt{ServingNode="None"})
    αντί να ανατεθεί στον "λιγότερο κακό" κόμβο· ο Πίνακας~\ref{tab:handovers}
    επιβεβαιώνει ότι το φαινόμενο της Ενότητας~\ref{subsec:results-outage}
    εξαλείφθηκε.
  \item \textbf{Πραγματική πολυσυνδεσιμότητα} (εξίσωση~\eqref{eq:node-selection},
    Ενότητα~\ref{sec:meth-node-selection}): το ίδιο κατώφλι
    $\text{SNR}_{\min}$ εφαρμόζεται πλέον \emph{ανεξάρτητα} στον
    καλύτερο BS και στον δορυφόρο, με τον χρήστη να συνδέεται σε
    όποιον/όποιους το ξεπερνούν -- ταυτόχρονα και στους δύο
    (\texttt{ServingType="DualConnectivity"}) όποτε το ξεπερνούν και οι
    δύο, με αθροιστική χωρητικότητα/ισχύ στις δύο ενεργές ζεύξεις
    (εξισώσεις~\eqref{eq:shannon-capacity},~\eqref{eq:energy-per-bit}).
    Πρόκειται για το τετρα-καταστασιακό μοντέλο σύνδεσης της
    αρχιτεκτονικής SS-SBS (Ενότητα~\ref{subsec:bg-architectures}), αντί
    για τον αρχικό, καθαρά single-connectivity κανόνα $\argmax$. Ένα
    απροσδόκητο εύρημα ήρθε μαζί με την αλλαγή: με τη γεωμετρία
    κάλυψης του σεναρίου αναφοράς, ένας χρήστης με χρησιμοποιήσιμο BS
    έχει σχεδόν πάντα \emph{και} χρησιμοποιήσιμο δορυφόρο, οπότε η
    κατηγορία Terrestrial-only εξαφανίστηκε εντελώς από το dataset της
    Ενότητας~\ref{sec:meth-montecarlo} -- κάτι που άλλαξε αισθητά και
    την ερμηνεία των πειραμάτων μηχανικής μάθησης
    (Ενότητα~\ref{sec:results-ml-realistic}), συζητείται στην επόμενη
    ενότητα.
  \item \textbf{Υστέρηση (hysteresis) και time-to-trigger} στην
    ενεργοποίηση/απενεργοποίηση ζεύξης
    (εξίσωση~\eqref{eq:hysteresis}, Ενότητα~\ref{subsec:meth-hysteresis}):
    μηχανή κατά το πρότυπο Event~A3 του TS~38.331~\cite{ts38331} πάνω
    στο ήδη υπάρχον κατώφλι $\text{SNR}_{\min}$, με περιθώριο
    $\Delta_{\text{hyst}}=2$~dB και $N_{\text{TTT}}=1$ βήμα στο σενάριο
    αναφοράς. Το ίδιο το σενάριο αναφοράς αποδείχθηκε "ήσυχο" (καθόλου
    ταλάντωση κοντά στο κατώφλι, αφού το επίγειο κανάλι παγώνει για
    ακίνητους χρήστες και το δορυφορικό είναι ήδη ντετερμινιστικό), οπότε
    η επικύρωση έγινε σε ένα ειδικά κατασκευασμένο σενάριο πίεσης
    (\texttt{hysteresisStressTest.m}, Ενότητα~\ref{subsec:results-hysteresis}):
    $67\%$ μείωση μεταβάσεων κατάστασης ($12\rightarrow4$ σε $80$
    βήματα), με εμφανή καταστολή ping-pong.
  \item \textbf{Joint/fairness-aware εξισορρόπηση φορτίου, σε δύο
    αλγορίθμους} (Ενότητα~\ref{subsec:meth-fairness}). \emph{PerBs}
    (εξίσωση~\eqref{eq:fairness}): όταν ένας BS υπερφορτώνεται, οι ήδη
    DualConnectivity χρήστες του (με δορυφορική εναλλακτική)
    αποσυνδέονται προτεραιοποιημένα από αυτόν -- ξεκινώντας από τον
    χαμηλότερο $\text{SNR}_{\text{BS}}$ -- αντί οι Terrestrial-only
    χρήστες (χωρίς εναλλακτική) να συνεχίζουν να μοιράζονται εξίσου το
    συρρικνωμένο εύρος ζώνης μαζί τους. Επικυρωμένος σε ειδικά
    κατασκευασμένο σενάριο συμφόρησης (\texttt{fairnessDemo.m},
    Ενότητα~\ref{subsec:results-fairness}): $+33\%$ μέση/ελάχιστη
    χωρητικότητα για την προστατευόμενη ομάδα, με ειλικρινά αναφερόμενο
    (όχι αποσιωπημένο) κόστος για την ομάδα με εναλλακτική. \emph{Joint}
    (εξίσωση~\eqref{eq:fairness-joint}), προστέθηκε στη συνέχεια πάνω στο
    ίδιο μηχανισμό: αντί να περιορίζεται σε έναν BS τη φορά, εντοπίζει
    τον πλέον υπερφορτωμένο BS σε όλο το δίκτυο και προτιμά lateral
    handover σε γείτονα BS με ελεύθερη χωρητικότητα αντί για πτώση σε
    δορυφόρο-μόνο -- κάτι που μπορεί να ωφελήσει ενεργά ακόμα και
    Terrestrial-only χρήστες, όχι μόνο να τους προστατέψει. Επικυρωμένος
    σε δεύτερο σενάριο δύο επικαλυπτόμενων BS
    (\texttt{jointFairnessDemo.m}, Ενότητα~\ref{subsec:results-fairness}):
    σε μια Terrestrial-only παραλλαγή το PerBs δεν μετακινεί κανέναν ενώ
    το Joint αναδιανέμει πλήρως το φορτίο· σε μια μεικτή DualConnectivity
    παραλλαγή κάθε μετακινημένος χρήστης βελτιώνεται σημαντικά
    ($1.6$-$4.3\times$), αλλά η μέση χωρητικότητα του πληθυσμού μπορεί να
    \emph{μειωθεί}, ένα ειλικρινά αναφερόμενο όριο του greedy
    heuristic -- ένα πραγματικό βήμα προς joint συντονισμό
    μεταξύ κόμβων, όχι όμως ακόμα πλήρης joint βελτιστοποίηση πολλαπλών
    κόμβων (βλ. Ενότητα~\ref{sec:disc-limitations} για τι παραμένει
    ανοιχτό).
\end{itemize}

\subsection{Η επίδραση της πολυσυνδεσιμότητας στα πειράματα μηχανικής μάθησης}
\label{subsec:disc-mc-ml-impact}

Πριν την προσθήκη της πολυσυνδεσιμότητας, τα τρία πειράματα μηχανικής
μάθησης της Ενότητας~\ref{sec:results-ml-realistic} σχημάτιζαν μια καθαρή,
μονότονη σκάλα ρεαλισμού: ακριβές SNR $\gg$ θορυβώδες SNR $>$ μόνο
γεωμετρία, με το Random Forest να πέφτει από σχεδόν τέλειο σε $89.5\%$ και
τελικά $87.2\%$ Accuracy. Μετά την προσθήκη της, η σκάλα αυτή είχε σχεδόν
εξαφανιστεί ($99.72\% \rightarrow 99.15\% \rightarrow 99.15\%$). Η αιτία
δεν ήταν μεθοδολογικό σφάλμα -- οι μετρικές ελέγχθηκαν ρητά ώστε να μη
συγχέουν έλλειψη δεδομένων με ευκολία προβλήματος (macro F1/ROC-AUC
περιορισμένα σε κλάσεις όντως παρούσες στο test set) -- αλλά πραγματική
συνέπεια της γεννήτριας σεναρίων εκείνης της στιγμής: η κατηγορία
Terrestrial-only είχε εξαφανιστεί εντελώς (βλ. παραπάνω), το τρι-κλασικό
πρόβλημα ανάγονταν ουσιαστικά σε DualConnectivity-έναντι-Satellite-only,
και η "κοντά"/"μακριά" τοποθέτηση χρηστών ήταν αυστηρά δίτιμη ($1$~km
πάντα-χρησιμοποιήσιμο έναντι $50$-$150$~km πάντα-εκτός εμβέλειας) --
καμία γνήσια οριακή περίπτωση δεν υπήρχε πουθενά στο dataset.

Η γεννήτρια Monte Carlo (\texttt{monteCarloDriver.m}) διορθώθηκε στη
συνέχεια ακριβώς σε αυτά τα δύο σημεία -- ευρύτερη τυχαιοποίηση του
υποδορυφορικού σημείου (ώστε ορισμένα σενάρια να μην έχουν καθόλου
δορυφορική κάλυψη) και μια ακτίνα "κοντινών" χρηστών $[0.1,4.0]$~km που
διασχίζει γνήσια το πραγματικό όριο χρησιμοποιήσιμου SNR -- και η σκάλα
ρεαλισμού \emph{επανεμφανίστηκε}: ακριβές SNR ($100.00\%$) $\gg$
μόνο-γεωμετρία ($93.93\%$) $\approx$ θορυβώδες SNR ($93.12\%$) για το
Random Forest, με τη Λογιστική Παλινδρόμηση να δείχνει την αναμενόμενη
σειρά καθαρότερα ($96.76\% > 94.74\% > 90.69\%$). Το ενδιάμεσο
"οροπέδιο" δεν ήταν συνεπώς ένα εγγενές χαρακτηριστικό της
πολυσυνδεσιμότητας αυτής καθαυτής, αλλά ένα τεχνούργημα της
συγκεκριμένης γεννήτριας σεναρίων που τη συνόδευε -- ένα εύρημα από μόνο
του, για το πόσο ευαίσθητη μπορεί να είναι η φαινομενική "δυσκολία" ενός
προβλήματος ταξινόμησης στις λεπτομέρειες του τρόπου παραγωγής των
δεδομένων εκπαίδευσης, ανεξάρτητα από την ίδια την αρχιτεκτονική
σύνδεσης. Η αναθεωρημένη τυπική απόκλιση θορύβου του πειράματος
θορυβώδους SNR (πλέον $\sigma_{SF}$ NLOS του 3GPP TR~38.901
Πίνακα 7.4.1-1 ανά \texttt{ScenarioType}, αντί για μια μετρημένη αλλά μη
αντιπροσωπευτική τιμή από παλαιότερη τοπολογία) εξηγείται στην
Ενότητα~\ref{sec:results-ml-realistic}.

Οι υπόλοιπες κατευθύνσεις παραμένουν ανοιχτές:

\begin{enumerate}
  \item \textbf{Πραγματικός ορβιτογράφος.} Αντικατάσταση του απλοποιημένου
    μοντέλου ίχνους εδάφους (Ενότητα~\ref{sec:meth-temporal}) με πλήρη
    διάδοση τροχιάς (π.χ. SGP4 πάνω σε πραγματικά two-line elements ενός
    LEO αστερισμού), ώστε οι χρονικές προσομοιώσεις να αντιστοιχούν σε
    πραγματικά, επαληθεύσιμα σενάρια διέλευσης.
  \item \textbf{Ιονοσφαιρική σπινθηρίδα και πλήρες μοτίβο κεραίας.} Η
    ατμοσφαιρική απόσβεση αερίων και το κέρδος κεραίας στο επίγειο σκέλος
    προστέθηκαν ήδη (βλ. παραπάνω)· παραμένουν ανοιχτά η ιονοσφαιρική
    σπινθηρίδα (TR~38.811 §6.6.6, χρειάζεται κλιματολογικό μοντέλο) και
    η αντικατάσταση του ενιαίου, βαθμωτού κέρδους κεραίας και στις δύο
    πλευρές με πλήρες, γωνιακά εξαρτώμενο μοτίβο ακτινοβολίας/πίνακα
    κεραιών (TR~38.901 §7.3), ώστε η ασυμμετρία διακύμανσης που
    παρατηρήθηκε στην Ενότητα~\ref{sec:results-repeated} να μην είναι
    απλώς τεχνούργημα ελλιπούς μοντελοποίησης.
  \item \textbf{Πλήρης joint βελτιστοποίηση πολλαπλών κόμβων και πιο
    σύνθετος διαμοιρασμός πόρων.} Ο αλγόριθμος Joint της
    Ενότητας~\ref{subsec:meth-fairness} συντονίζει ήδη πολλαπλούς BS
    (worst-BS-first, lateral handover αντί για πτώση σε δορυφόρο-μόνο
    όταν υπάρχει επικάλυψη κάλυψης) -- το κενό που απομένει είναι πλέον
    πιο συγκεκριμένο, όχι η πλήρης απουσία συντονισμού: χρειάζεται μια
    πραγματική αντικειμενική συνάρτηση πάνω στη συνολική/βεβαρημένη
    χωρητικότητα του δικτύου (π.χ. proportional fairness πάνω στο
    throughput) αντί για την τρέχουσα greedy επιλογή "προορισμός με τη
    μεγαλύτερη ελεύθερη χωρητικότητα", η οποία -- όπως έδειξε το
    \texttt{jointFairnessDemo.m} (Ενότητα~\ref{subsec:results-fairness}) --
    μπορεί να μειώσει τη μέση χωρητικότητα του πληθυσμού αραιώνοντας
    έναν προηγουμένως ελαφρά φορτωμένο γείτονα BS. Ταυτόχρονη
    εξισορρόπηση BS \emph{και} δορυφόρου μαζί (όχι μόνο ο δορυφόρος ως
    fallback προορισμός) και πιο σύνθετα σχήματα διαμοιρασμού (π.χ.
    Mode~2/3 της SS-SBS ή η αρχιτεκτονική MS-MBS,
    Ενότητα~\ref{subsec:bg-architectures}, με κοινή διαχείριση
    φάσματος/ισχύος μεταξύ πολλαπλών δορυφόρων ή σταθμών βάσης)
    παραμένουν επίσης ανοιχτά.
  \item \textbf{Πέρα από ταξινόμηση \texttt{ServingType}.} Η γεννήτρια
    Monte Carlo πλέον παράγει γνήσιες οριακές περιπτώσεις και τις τρεις
    κλάσεις με πραγματική συχνότητα (βλ. Ενότητα~\ref{subsec:disc-mc-ml-impact}
    παραπάνω), οπότε το πιο επείγον σκέλος αυτού του σημείου έχει ήδη
    αντιμετωπιστεί -- ο ίδιος ο στόχος πρόβλεψης παραμένει όμως η στατική
    ταξινόμηση \texttt{ServingType}, μια ουσιωδώς απλή ερώτηση σε σχέση
    με ό,τι θα χρειαζόταν πραγματικά ένα σύστημα διαχείρισης
    συνδεσιμότητας. Ουσιωδώς πιο ρεαλιστικά προβλήματα παραμένουν
    ανοιχτά: πρόβλεψη handover/SN event \emph{πριν} συμβεί, δηλαδή με
    χρονική προήγηση πάνω στα δεδομένα της χρονικής προσομοίωσης της
    Ενότητας~\ref{sec:results-temporal}, μετάβαση από ταξινόμηση σε
    regression/ranking στόχους (\texttt{Capacity\_Mbps},
    \texttt{EnergyPerBit\_uJ}), ή ένας joint/fairness-aware στόχος πάνω
    στο μηχανισμό της Ενότητας~\ref{subsec:meth-fairness} (π.χ. πρόβλεψη
    αν ένας χρήστης θα μετακινηθεί λόγω φόρτου). Κανένα από αυτά δεν
    βελτιστοποιεί η τρέχουσα SNR-greedy baseline.
\end{enumerate}

Τα σημεία 1-3 παραπάνω απαιτούν πιο εκτεταμένη επέκταση του μοντέλου
link-budget/κίνησης δορυφόρου/βελτιστοποίησης, ενώ το σημείο 4 (μηχανική
μάθηση σε μη-τετριμμένο στόχο) είναι το φυσικό Μέρος~2 της εργασίας και
θα ακολουθήσει αφού σταθεροποιηθεί περαιτέρω το μοντέλο προσομοίωσης
(Μέρος~1). Τόσο η υστέρηση/TTT όσο και η εξισορρόπηση φορτίου (PerBs \emph{και}
Joint), που ως πρόσφατα ήταν τα σημεία με την υψηλότερη προτεραιότητα σε
αυτή τη λίστα, υλοποιήθηκαν ήδη (βλ. παραπάνω) -- η πλήρης joint
βελτιστοποίηση πολλαπλών κόμβων (σημείο 3) παραμένει η φυσική επόμενη
προτεραιότητα σε επίπεδο κανόνα επιλογής κόμβου, πλέον όμως πιο
συγκεκριμένη: όχι η εισαγωγή συντονισμού μεταξύ BS από το μηδέν (αυτό
έγινε ήδη με το Joint), αλλά η αντικατάσταση του τρέχοντος greedy
heuristic με μια πραγματική αντικειμενική συνάρτηση.
